%% LyX 2.5.1 created this file.  For more info, see https://www.lyx.org/.
%% Do not edit unless you really know what you are doing.
\documentclass[brazil]{article}
\usepackage[T1]{fontenc}
\usepackage{textcomp}
\usepackage[utf8]{inputenc}
\usepackage{babel}
\usepackage{amsmath}
\usepackage{amssymb}
\usepackage{geometry}
\geometry{verbose,tmargin=3cm,bmargin=3cm,lmargin=3cm,rmargin=2.5cm}
\usepackage[bookmarks=false,
 breaklinks=false,pdfborder={0 0 1},backref=false,colorlinks=false]
 {hyperref}

\makeatletter
%%%%%%%%%%%%%%%%%%%%%%%%%%%%%% User specified LaTeX commands.


\makeatother

\begin{document}
\title{A função de onda}
\maketitle

\subsection{Equação de Schrödinger}

A equação de Schrödinger desempenha um papel análogo na mecânica quântica ao da segunda lei de Newton na mecânica clássica:

\begin{equation}
i\hbar\frac{\partial\Psi}{\partial t}=-\frac{\hbar^{2}}{2m}\frac{\partial^{2}\Psi}{\partial x^{2}}+V\Psi\label{eq:Um}
\end{equation}

Onde $\hbar$ é a constante de Planck dividida por $2\pi$. Ou seja, se temos uma partícula de massa $m$ se movendo ao longo do eixo $x$ e queremos obter informações sobre a partícula em qualquer instante $t$, a partir de uma condição inicial adequada (tipicamente $\Psi\left(x,0\right)$), a equação \ref{eq:Um} determina a função de onda $\Psi\left(x,t\right)$ para todo tempo futuro, de forma análoga a como a segunda lei de Newton determina a posição $x\left(t\right)$ para todo instante futuro.

Vale destacar que a função de onda não é algo que pode ser medido diretamente, o que fazemos é usá-la para obter grandezas que podem ser observadas diretamente. Por exemplo, a interpretação estatística de Born nos fornece a probabilidade de encontrar a partícula entre os pontos $a$ e $b$ no instante $t$:

\begin{equation}
P=\int^{b}_{a}\left|\Psi\left(x,t\right)\right|^{2}dx
\end{equation}

Como $\left|\Psi\left(x,t\right)\right|^{2}$ representa uma densidade de probabilidade no espaço, ela deve ter unidade de $m^{-1}$ em uma dimensão. Consequentemente, a função de onda $\Psi\left(\boldsymbol{x},t\right)$ tem unidade de $m^{-1/2}$. De forma mais geral, em $n$ dimensões, a função de onda tem unidade de $m^{-n/2}$. Como a partícula deve estar em algum lugar, então:

\begin{equation}
\int^{+\infty}_{-\infty}\left|\Psi\left(x,t\right)\right|^{2}dx=1\label{eq:2}
\end{equation}

Retornando à equação de Schrödinger: se $\Psi$ é uma solução, então $A\Psi$ também será uma solução para qualquer constante $A$. Podemos escolher o valor de $A$ de modo que a condição de normalização seja satisfeita. Se a solução da equação de Schrödinger diverge (isto é, tende ao infinito) ou é nula em todo o espaço, ela não é normalizável e, portanto, não representa um estado físico conhecido e deve ser descartada. Vale destacar que, uma vez normalizada, a função de onda $\Psi$ permanecerá normalizada ao longo do tempo. Isso reflete a conservação da probabilidade na mecânica quântica.

\begin{equation}
\frac{d}{dt}\int^{+\infty}_{-\infty}\left|\Psi\left(x,t\right)\right|^{2}dx=\int^{+\infty}_{-\infty}\frac{\partial}{\partial t}\left|\Psi\left(x,t\right)\right|^{2}dx
\end{equation}

Aqui, fizemos a troca de derivada total por derivada parcial dentro da integral porque a integral resulta em uma função que depende apenas do tempo $t$, enquanto o integrando $|\Psi(x,t)|^{2}$ depende tanto de $x$ quanto de $t$.

\begin{equation}
\frac{\partial}{\partial t}\left|\Psi\right|^{2}=\frac{\partial}{\partial t}\left|\Psi\right|^{2}=\frac{\partial}{\partial t}\left(\sqrt{\Psi^{*}\Psi}\right)^{2}=\frac{\partial}{\partial t}\left(\Psi^{*}\Psi\right)=\Psi^{*}\frac{\partial\Psi}{\partial t}+\frac{\partial\Psi^{*}}{\partial t}\Psi
\end{equation}

A equação de Schrödinger é:

\begin{equation}
\frac{\partial\Psi}{\partial t}=\frac{i\hbar}{2m}\frac{\partial^{2}\Psi}{\partial x^{2}}-\frac{i}{\hbar}V\Psi
\end{equation}

E tomando o complexo conjugado, assumindo $V$ real:

\begin{equation}
\frac{\partial\Psi^{*}}{\partial t}=-\frac{i\hbar}{2m}\frac{\partial^{2}\Psi^{*}}{\partial x^{2}}+\frac{i}{\hbar}V\Psi^{*}
\end{equation}

Então:

\begin{align}
\frac{\partial}{\partial t}\left|\Psi\right|^{2} & =\Psi^{*}\frac{\partial\Psi}{\partial t}+\frac{\partial\Psi^{*}}{\partial t}\Psi\\
 & =\Psi^{*}\left(\frac{i\hbar}{2m}\frac{\partial^{2}\Psi}{\partial x^{2}}-\frac{i}{\hbar}V\Psi\right)+\left(-\frac{i\hbar}{2m}\frac{\partial^{2}\Psi^{*}}{\partial x^{2}}+\frac{i}{\hbar}V\Psi^{*}\right)\Psi\\
 & =\Psi^{*}\frac{i\hbar}{2m}\frac{\partial^{2}\Psi}{\partial x^{2}}-\frac{i\hbar}{2m}\frac{\partial^{2}\Psi^{*}}{\partial x^{2}}\Psi\\
 & =\frac{i\hbar}{2m}\left(\Psi^{*}\frac{\partial^{2}\Psi}{\partial x^{2}}-\frac{\partial^{2}\Psi^{*}}{\partial x^{2}}\Psi\right)\label{eq:nl}
\end{align}

Então:

\begin{align}
\frac{d}{dt}\int^{+\infty}_{-\infty}\left|\Psi\left(x,t\right)\right|^{2}dx & =\int^{+\infty}_{-\infty}\frac{\partial}{\partial t}\left|\Psi\left(x,t\right)\right|^{2}dx\\
 & =\int^{+\infty}_{-\infty}\frac{i\hbar}{2m}\left(\Psi^{*}\frac{\partial^{2}\Psi}{\partial x^{2}}-\frac{\partial^{2}\Psi^{*}}{\partial x^{2}}\Psi\right)dx\\
 & =\frac{i\hbar}{2m}\int^{+\infty}_{-\infty}\left(\Psi^{*}\frac{\partial^{2}\Psi}{\partial x^{2}}-\frac{\partial^{2}\Psi^{*}}{\partial x^{2}}\Psi\right)dx
\end{align}

Olhando o primeiro termo, vamos aplicar integral por partes (\href{https://proofwiki.org/wiki/Integration_by_Parts}{prova}):

\begin{equation}
\int^{b}_{a}f\left(x\right)g'\left(x\right)dx=\left.f\left(x\right)g\left(x\right)\right|^{b}_{a}-\int^{b}_{a}g\left(x\right)f'\left(x\right)dx
\end{equation}

Onde também temos que $u=f\left(x\right)$ e $dv=f'\left(x\right)dx$. Então sendo $f=\Psi^{*}$ e $g=\Psi'$

\begin{equation}
\int^{+\infty}_{-\infty}\Psi\left(x,t\right)^{*}\frac{\partial^{2}\Psi\left(x,t\right)}{\partial x^{2}}dx=\left.\Psi^{*}\Psi'\right|^{+\infty}_{-\infty}-\int^{+\infty}_{-\infty}\Psi'\Psi^{*}{}'dx
\end{equation}

De modo análogo temos pra segunda integral $f=\Psi$ e $g=\Psi^{*}{}'$

\begin{equation}
\int^{+\infty}_{-\infty}\Psi\frac{\partial^{2}\Psi^{*}}{\partial x^{2}}dx=\left.\Psi\Psi^{*}{}'\right|^{+\infty}_{-\infty}-\int^{+\infty}_{-\infty}\Psi^{*}{}'\Psi'dx
\end{equation}

Então:

\begin{align}
\frac{d}{dt}\int^{+\infty}_{-\infty}\left|\Psi\left(x,t\right)\right|^{2}dx & =\frac{i\hbar}{2m}\int^{+\infty}_{-\infty}\left(\Psi^{*}\frac{\partial^{2}\Psi}{\partial x^{2}}-\frac{\partial^{2}\Psi^{*}}{\partial x^{2}}\Psi\right)dx\\
 & =\frac{i\hbar}{2m}\left[\left.\Psi^{*}\Psi'\right|^{+\infty}_{-\infty}-\int^{+\infty}_{-\infty}\Psi'\Psi^{*}{}'dx-\left.\Psi\Psi^{*}{}'\right|^{+\infty}_{-\infty}+\int^{+\infty}_{-\infty}\Psi^{*}{}'\Psi'dx\right]\\
 & =\frac{i\hbar}{2m}\left[\left.\Psi^{*}\Psi'\right|^{+\infty}_{-\infty}-\left.\Psi\Psi^{*}{}'\right|^{+\infty}_{-\infty}\right]\\
 & =\frac{i\hbar}{2m}\left.\left(\Psi^{*}\Psi'-\Psi\Psi^{*}{}'\right)\right|^{+\infty}_{-\infty}
\end{align}

Para que um resultado seja fisicamente válido, a função de onda e suas derivadas devem tender a zero no infinito.\footnote{Existem construções matemáticas de funções de onda que podem ser normalizadas, mesmo sem tender exatamente a zero no infinito. Contudo, segundo Griffiths, tais funções não aparecem na física, pois uma função de onda física deve refletir a probabilidade de encontrar uma partícula em alguma região finita do espaço. Funções que não decaem no infinito não correspondem a estados físicos observáveis.}, então:

\begin{align}
\frac{d}{dt}\int^{+\infty}_{-\infty}\left|\Psi\left(x,t\right)\right|^{2}dx & =0
\end{align}

Ou seja, como o valor da integral é constante ao longo do tempo, se a função de onda estiver normalizada no instante $t=0$, ela permanecerá normalizada para todos os instantes seguintes.

\textbf{Exemplo: Considere a função de onda}

\begin{equation}
\Psi\left(x,t\right)=Ae^{-\lambda\left|x\right|}e^{-i\omega t}
\end{equation}

\textbf{Onde $A$, $\lambda$ e $\omega$ são constantes reais e positivas. Vamos calcular:}

\textbf{a) Normalização de $\Psi$.}

\textbf{b) Determinar os valores esperados de $x$ e $x^{2}$.}

\textbf{c) Encontrar o desvio padrão de $x$ e calcular a probabilidade da partícula ser encontrada fora do intervalo entre $\left\langle x\right\rangle -\sigma$ e $\left\langle x\right\rangle +\sigma$.}

Sendo $x$ um número real então:

\begin{align}
\int^{+\infty}_{-\infty}\left|\Psi\left(x,t\right)\right|^{2}dx & =1\\
\int^{+\infty}_{-\infty}\Psi\Psi^{*}dx & =\\
\int^{+\infty}_{-\infty}\left(Ae^{-\lambda\left|x\right|}e^{-i\omega t}\right)\left(Ae^{-\lambda\left|x\right|}e^{i\omega t}\right)dx & =\\
A^{2}\int^{+\infty}_{-\infty}e^{-\lambda2\left|x\right|}dx & =\\
A^{2}\left(\int^{0}_{-\infty}e^{-\lambda2\left(-x\right)}dx+\int^{+\infty}_{0}e^{-\lambda2x}dx\right) & =\\
A^{2}\left(\frac{1}{2\lambda}\int^{0}_{-\infty}e^{u}du+\frac{1}{-2\lambda}\int^{+\infty}_{0}e^{v}dv\right) & =
\end{align}

Chamo atenção para a 5ª linha onde dividimos a integral para respeitar o módulo, usando a propriedade de aditividade da integral\footnote{$\int^{c}_{a}f\left(x\right)dx=\int^{b}_{a}f\left(x\right)dx+\int^{c}_{b}f\left(x\right)dx$ onde $a\leq b\leq c$ (\href{https://www.math.ucdavis.edu/~hunter/m125b/ch1.pdf}{prova}).}. Onde no primeiro caso temos $u=2\lambda x$ e no segundo caso $v=-2\lambda x$

\begin{align}
A^{2}\left(\frac{1}{2\lambda}\left.e^{2\lambda x}\right|^{0}_{-\infty}+\frac{1}{-2\lambda}\left.e^{-2\lambda x}\right|^{+\infty}_{0}\right) & =1\\
A^{2}\left(\frac{\left(e^{0}-\frac{1}{e^{\infty}}\right)}{2\lambda}+\frac{\left(\frac{1}{e^{\infty}}-e^{0}\right)}{-2\lambda}\right) & =\\
A^{2}\left(\frac{1}{2\lambda}+\frac{1}{2\lambda}\right) & =\\
\frac{A^{2}}{\lambda} & =\\
A & =\sqrt{\lambda}
\end{align}

O valor esperado para $x$ é então:

\begin{align}
\int^{+\infty}_{-\infty}x\left|\Psi\left(x,t\right)\right|^{2}dx & =A^{2}\left(\int^{0}_{-\infty}xe^{2\lambda x}dx+\int^{+\infty}_{0}xe^{-2\lambda x}dx\right)\\
 & =\lambda\left(\int^{0}_{-\infty}\frac{1}{2\lambda}\frac{u}{2\lambda}e^{u}du+\int^{+\infty}_{0}\frac{1}{2\lambda}\frac{u}{2\lambda}e^{-u}du\right)\\
 & =\frac{1}{4\lambda}\left(\int^{0}_{-\infty}ue^{u}du+\int^{+\infty}_{0}ue^{-u}du\right)\label{eq:5}
\end{align}

Onde usamos novamente que $u=2\lambda x$. Poderíamos usar apenas o fato de que o integrando ser ímpar e termos limites simétricos\footnote{Funções pares são funções que $f\left(-x\right)=f\left(x\right),$temos então $\int^{a}_{-a}f\left(x\right)dx=2\int^{a}_{0}f\left(x\right)dx$. Funções ímpares são funções que $f\left(-x\right)=-f\left(x\right)$, temos então $\int^{a}_{-a}f\left(x\right)dx=0$ (\href{https://proofwiki.org/wiki/Definite_Integral_of_Odd_Function}{prova}e \href{https://proofwiki.org/wiki/Definite_Integral_of_Even_Function}{prova}).}. Mas queremos resolver por completo. Sendo $f\left(u\right)=u$ e $g\left(u\right)=e^{u}$

\begin{align}
\int^{b}_{a}f\left(u\right)g'\left(u\right)du & =\left.f\left(u\right)g\left(u\right)\right|^{b}_{a}-\int^{b}_{a}g\left(u\right)f'\left(u\right)du\\
\int^{0}_{-\infty}ue^{u}du & =\left.ue^{u}\right|^{0}_{-\infty}-\int^{9}_{-\infty}e^{u}du\\
 & =\lim_{u\rightarrow-\infty}\left(ue^{u}\right)-\int^{9}_{-\infty}e^{u}du\\
 & =\lim_{u\rightarrow\infty}\left(-\frac{u}{e^{u}}\right)-\left(e^{0}-\frac{1}{e^{\infty}}\right)\\
 & =\lim_{u\rightarrow\infty}\left(-\frac{u}{e^{u}}\right)-1
\end{align}

Como temos uma indeterminação do tipo $-\infty/\infty$ podemos aplicar a regra de L'hopital\footnote{$\lim_{x\rightarrow a}\frac{f\left(x\right)}{g\left(x\right)}=\lim_{x\rightarrow a}\frac{f'\left(x\right)}{g'\left(x\right)},$quando temos uma indeterminação do tipo $0/0$ ou $\pm\infty/\pm\infty$ (\href{https://proofwiki.org/wiki/L'H\%5C\%2525C3\%5C\%2525B4pital's_Rule}{prova}).}:

\begin{align}
\int^{0}_{-\infty}ue^{u}du & =\lim_{u\rightarrow\infty}\left(-\frac{du/dt}{de^{u}/dt}\right)-1\\
 & =\lim_{u\rightarrow\infty}\left(-\frac{du/dt}{de^{u}/dt}\right)-1\\
 & =\lim_{u\rightarrow\infty}\left(-\frac{1}{e^{u}}\right)-1\\
 & =-1
 \label{eq:9}
\end{align}

E equivalente:

\begin{align}
\int^{+\infty}_{0}ue^{-u}du & =\left.-ue^{-u}\right|^{+\infty}_{0}+\int^{+\infty}_{0}e^{-u}du\\
 & =-\lim_{u\rightarrow\infty}\left(ue^{-u}\right)-\left(\frac{1}{e^{\infty}}-1\right)\\
 & =-\lim_{u\rightarrow\infty}\left(\frac{u}{e^{u}}\right)+1\\
 & =1\label{eq:13}
\end{align}

Este é o mesmo limite que calculamos anteriormente aplicando a L'hopital e temos $\lim_{u\rightarrow\infty}\left(u/e^{u}\right)=0$. Então retomando a equação \ref{eq:5} temos:

\begin{align}
\int^{+\infty}_{-\infty}x\left|\Psi\left(x,t\right)\right|^{2}dx & =\frac{1}{4\lambda}\left(\int^{0}_{-\infty}ue^{u}du+\int^{+\infty}_{0}ue^{-u}du\right)\\
 & =\frac{1}{4\lambda}\left(-1+1\right)\\
 & =0
\end{align}

Temos algo análogo para o valor de esperado de $x^{2}$ onde vamos usar a mesma substituição $u=2\lambda x$

\begin{align}
\int^{+\infty}_{-\infty}x^{2}\left|\Psi\left(x,t\right)\right|^{2}dx & =A^{2}\left(\int^{0}_{-\infty}x^{2}e^{2\lambda x}dx+\int^{+\infty}_{0}x^{2}e^{-2\lambda x}dx\right)\\
 & =\lambda\left(\int^{0}_{-\infty}\frac{1}{2\lambda}\left(\frac{u}{2\lambda}\right)^{2}e^{u}du+\int^{+\infty}_{0}\frac{1}{2\lambda}\left(\frac{u}{2\lambda}\right)^{2}e^{-u}du\right)\\
 & =\frac{1}{8\lambda^{2}}\left(\int^{0}_{-\infty}u^{2}e^{u}du+\int^{+\infty}_{0}u^{2}e^{-u}du\right)\label{eq:16}
\end{align}

Aplicando novamente a integração por partes mas sendo agora $f\left(u\right)=u^{2}$e $g\left(u\right)=e^{u}$ para a primeira integral:

\begin{align}
\int^{b}_{a}f\left(x\right)g'\left(x\right)dx & =\left.f\left(x\right)g\left(x\right)\right|^{b}_{a}-\int^{b}_{a}g\left(x\right)f'\left(x\right)dx\\
\int^{0}_{-\infty}u^{2}e^{u}du & =\left.u^{2}e^{u}\right|^{0}_{-\infty}-2\int^{0}_{-\infty}ue^{u}du\\
 & =\lim_{u\rightarrow\infty}\frac{u^{2}}{e^{u}}-2\int^{0}_{-\infty}ue^{u}due^{u}
\end{align}

O resultado desta integral já foi encontrado na equação \ref{eq:9} e é $-1$. O limite podemos usar duas vezes l'hopital, de forma que obtemos apenas:

\begin{equation}
\int^{0}_{-\infty}u^{2}e^{u}du=2\label{eq:17}
\end{equation}

Para a segunda integral temos:

\begin{align}
\int^{b}_{a}f\left(x\right)g'\left(x\right)dx & =\left.f\left(x\right)g\left(x\right)\right|^{b}_{a}-\int^{b}_{a}g\left(x\right)f'\left(x\right)dx\\
\int^{+\infty}_{0}u^{2}e^{-u}du & =-\left.u^{2}e^{-u}\right|^{\infty}_{0}+2\int^{+\infty}_{0}ue^{-u}du\\
 & =-\lim_{u\rightarrow\infty}\frac{u^{2}}{e^{u}}+2\int^{0}_{-\infty}ue^{-u}du
\end{align}

De forma análoga podemos aplicar duas vezes l'hopital para resolver o limite e o resultado dessa integral foi obtido em \ref{eq:9} e vale $1$. Temos então:

\begin{equation}
\int^{+\infty}_{0}u^{2}e^{-u}du=2\label{eq:18}
\end{equation}

Agora combinando os resultados \ref{eq:17} e \ref{eq:18} em \ref{eq:16}:

\begin{align}
\int^{+\infty}_{-\infty}x^{2}\left|\Psi\left(x,t\right)\right|^{2}dx & =\frac{1}{8\lambda^{2}}\left(\int^{0}_{-\infty}u^{2}e^{u}du+\int^{+\infty}_{0}u^{2}e^{-u}du\right)\\
 & =\frac{1}{8\lambda^{2}}\left(2+2\right)\\
 & =\frac{1}{2\lambda^{2}}
\end{align}

O desvio padrão de $x$ é dado então por:

\begin{align}
\sigma & =\sqrt{\left\langle x^{2}\right\rangle -\left\langle x\right\rangle ^{2}}
\end{align}

Como $\left\langle x\right\rangle =\int^{+\infty}_{-\infty}x\left|\Psi\left(x,t\right)\right|^{2}=0$ e $\left\langle x^{2}\right\rangle =\int^{+\infty}_{-\infty}x^{2}\left|\Psi\left(x,t\right)\right|^{2}=1/2\lambda^{2}$ então:

\begin{equation}
\sigma=\sqrt{\frac{1}{2\lambda^{2}}-0}=\frac{1}{\sqrt{2}\lambda}
\end{equation}

E por fim a probabilidade de não estar no intervalo entre $-\sigma$ e $+\sigma$ é dado por:

\begin{align}
\int^{--\sigma}_{-\infty}\left|\Psi\right|^{2}dx+\int^{+\infty}_{+\sigma}\left|\Psi\right|^{2}dx & =-\int^{--\infty}_{-\sigma}\left|\Psi\left(x,t\right)\right|^{2}dx+\int^{+\infty}_{+\sigma}\Psi\left(x,t\right)^{2}dx\\
 & =\int^{+\infty}_{+\sigma}\Psi\left(-u,t\right)^{2}du+\int^{+\infty}_{+\sigma}\Psi\left(x,t\right)^{2}dx
\end{align}

Onde usamos a propriedade sobre inverter os limites da integral\footnote{$\int^{b}_{a}f\left(x\right)dx=-\int^{a}_{b}f\left(x\right)dx$ (\href{https://proofwiki.org/wiki/Reversal_of_Limits_of_Definite_Integral}{prova}).}, e fizemos a substituição $u=-x,$ de forma que temos $du/dx=-1$ ou $du=-dx$ e os novos limites são 
\begin{equation}
\lim_{x\rightarrow-v}u=\lim_{x\rightarrow-v}-x=v
\end{equation}

No limite inferior $v=\sigma$ e no limite superior $v=-\infty$. No nosso caso como $\Psi\left(x,t\right)$ é uma função par, isto é $\Psi\left(-x,t\right)=\Psi\left(x,t\right)$, pois $x$ sempre aparece em módulo $\left|x\right|$ então:

\begin{align}
\int^{--\sigma}_{-\infty}\left|\Psi\right|^{2}dx+\int^{+\infty}_{+\sigma}\left|\Psi\right|^{2}dx & =\int^{+\infty}_{+\sigma}\Psi\left(u,t\right)^{2}du+\int^{+\infty}_{+\sigma}\Psi\left(x,t\right)^{2}dx\\
 & =2\int^{+\infty}_{+\sigma}\Psi\left(x,t\right)^{2}dx
\end{align}

Onde temos o segundo resultado pois $x$ e $v$ são apenas símbolo para uma variável, que nos dois casos são a mesma variável na função $\Psi$, e temos consequentemente a mesma integral. Resolvendo então, essa é a mesma integral indeterminada resolvida no começo do exemplo, com a vantagem de que $0<\sigma<\infty$, ou seja, os valores dentro dos limites da integral são sempre positivos de forma que dentro dessa faixa de valores podemos escrever que $\left|x\right|=x$.

\begin{align}
\int^{--\sigma}_{-\infty}\left|\Psi\right|^{2}dx+\int^{+\infty}_{+\sigma}\left|\Psi\right|^{2}dx & =2\int^{+\infty}_{+\sigma}\Psi\left(x,t\right)^{2}dx\\
 & =2A^{2}\int^{+\infty}_{+\sigma}e^{-\lambda2\left|x\right|}dx\\
 & =2\lambda\int^{+\infty}_{+\sigma}e^{-\lambda2x}dx\\
 & =\frac{2\lambda}{-2\lambda}\left[\frac{1}{e^{\lambda2x}}\right]^{+\infty}_{+\sigma}\\
 & =-\left[\frac{1}{e^{\infty}}-\frac{1}{e^{\lambda2\sigma}}\right]^{+\infty}_{+\sigma}\\
 & =\exp\left[-\lambda2\sigma\right]\\
 & =\exp\left[-\frac{2\lambda}{\sqrt{2}\lambda}\right]
\end{align}

Onde usamos o fato de que calculamos já $A$ e $\sigma.$ Temos então que a probabilidade é e$^{-\sqrt{2}}$ou $0.2431$, ou ainda $24,31\%$.

\subsection{Momentum}

O valor esperado $\left\langle x\right\rangle $ é a média das medidas feitas em uma coleção de sistemas idênticos, e não a média de medidas repetidas em um mesmo sistema. Porém, agora $\left\langle x\right\rangle $ varia com o tempo, e podemos querer saber quão rápido ele se move, de modo análogo ao que fazemos na mecânica clássica, onde a partir de $\ddot{x}$ investigamos outras propriedades do sistema. Temos então:

\begin{align}
\frac{d\left\langle x\right\rangle }{dt} & =\int x\frac{\partial}{\partial t}\left|\Psi\left(x,t\right)\right|^{2}
\end{align}

Onde estamos derivando apenas em relação ao que depende explicitamente do tempo, e aqui há algo que talvez nos pareça estranho. Diferentemente da mecânica clássica, em que temos $x(t)$, nesta abordagem matemática da mecânica quântica, $x$ não depende do tempo; não é uma função de posição. Ele é apenas um rótulo para as diferentes posições que uma partícula pode ocupar. Portanto, não há uma evolução temporal de $x$. Os rótulos $x$ para $\Psi(x,0)$ não evoluem para os rótulos de $\Psi(x,1)$, assim como as diferentes posições $x$ para $\Psi(x,0)$ são independentes entre si. Ou seja, para $x=1$, sempre teremos $x=1$; não existe uma função $x(t)$ que evolua tal que $x(0)=0$ e $x(1)=1$.

O que ocorre é que $\Psi(0,0)$, $\Psi(1,0)$, $\Psi(0,1)$ e $\Psi(1,1)$ são todos valores distintos da função $\Psi(x,t)$. É a função de onda que evolui no tempo, e nela podemos obter diferentes probabilidades de medir a partícula tanto em $x=0$ quanto em $x=1$, e tanto em $t=0$ quanto em $t=1$. Então usando o resultado \ref{eq:nl}:

\begin{equation}
\frac{\partial}{\partial t}\left|\Psi\right|^{2}=\frac{i\hbar}{2m}\left(\Psi^{*}\frac{\partial^{2}\Psi}{\partial x^{2}}-\frac{\partial^{2}\Psi^{*}}{\partial x^{2}}\Psi\right)
\end{equation}

Temos:

\begin{align}
\frac{d\left\langle x\right\rangle }{dt} & =\int x\frac{\partial}{\partial t}\left|\Psi\left(x,t\right)\right|^{2}\\
 & =\frac{i\hbar}{2m}\int x\left(\Psi^{*}\frac{\partial^{2}\Psi}{\partial x^{2}}-\frac{\partial^{2}\Psi^{*}}{\partial x^{2}}\Psi\right)dx
\end{align}

Onde usando a regra do produto:

\begin{equation}
\frac{d}{dx}\left(fg\right)=f\frac{dg}{dx}+\frac{df}{dx}g
\end{equation}

Então:

\begin{align}
\frac{\partial}{\partial x}\left(\Psi^{*}\frac{\partial\Psi}{\partial x}-\frac{\partial\Psi^{*}}{\partial x}\Psi\right) & =\frac{\partial}{\partial x}\left(\Psi^{*}\frac{\partial\Psi}{\partial x}\right)-\frac{\partial}{\partial x}\left(\frac{\partial\Psi^{*}}{\partial x}\Psi\right)\\
 & =\left[\frac{\partial\Psi^{*}}{\partial x}\frac{\partial\Psi}{\partial x}+\Psi^{*}\frac{\partial^{2}\Psi}{\partial x^{2}}\right]-\left[\frac{\partial^{2}\Psi^{*}}{\partial x^{2}}\Psi+\frac{\partial\Psi^{*}}{\partial x}\frac{\partial\Psi}{\partial x}\right]\\
 & \frac{\partial\Psi^{*}}{\partial x}\frac{\partial\Psi}{\partial x}+\Psi^{*}\frac{\partial^{2}\Psi}{\partial x^{2}}-\frac{\partial^{2}\Psi^{*}}{\partial x^{2}}\Psi
\end{align}

Então podemos reescrever a integral como:

\begin{equation}
\frac{d\left\langle x\right\rangle }{dt}=\frac{i\hbar}{2m}\int x\left[\frac{\partial}{\partial x}\left(\Psi^{*}\frac{\partial\Psi}{\partial x}-\frac{\partial\Psi^{*}}{\partial x}\Psi\right)\right]dx
\end{equation}

Por conveniência não estamos escrevendo os limites da integral $\int^{+\infty}_{-\infty}\rightarrow\int$ e nem as dependências da função de onda $\Psi\left(x,t\right)=\Psi$. Usando integral por partes então em cada uma das integrais escolhendo $f\left(x\right)=x$ e $g'\left(x\right)=\frac{\partial}{\partial x}\left(\Psi^{*}\frac{\partial\Psi}{\partial x}-\frac{\partial\Psi^{*}}{\partial x}\Psi\right)$, para a primeira integral, temos:

\begin{align}
\int^{b}_{a}f\left(x\right)g'\left(x\right)dx & =\left.f\left(x\right)g\left(x\right)\right|^{b}_{a}-\int^{b}_{a}g\left(x\right)f'\left(x\right)dx\\
\int x\left(\Psi^{*}\frac{\partial^{2}\Psi}{\partial x^{2}}\right) & =\left[x\left(\Psi^{*}\frac{\partial\Psi}{\partial x}-\frac{\partial\Psi^{*}}{\partial x}\Psi\right)\right]^{+\infty}_{-\infty}-\int\left(\Psi^{*}\frac{\partial\Psi}{\partial x}-\frac{\partial\Psi^{*}}{\partial x}\Psi\right)dx
\end{align}

Uma demonstração mais robusta matematicamente pode ser necessária futuramente, mas exigimos que $\Psi(x)$ (e suas derivadas) vá a zero quando $x\to\infty$ mais rápido do que o próprio $x$ vai para infinito --- ou seja, mais rápido do que $1/|x|\to0$ quando $x\to\infty$. Essa é uma exigência relacionada também às condições de normalização da função de onda. Apesar de não termos discutido o quão rápido a função de onda deve tender a zero no infinito, já havíamos mencionado essa exigência para que a função de onda consiga descrever um estado fisicamente possível. Aqui, estamos apenas assumindo que isso ocorre mais rápido do que $1/|x|$. Além disso, em sistemas típicos, a derivada de $\Psi$ se comporta de forma semelhante. É importante notar que aqui assumimos várias coisas sobre a função de onda para que esse resultado seja válido, de forma sucinta: 
\begin{itemize}
\item $\Psi$ e suas derivadas vão a zero quando $x\to\infty$ mais rápido que $1/|x|$. 
\end{itemize}
Podemos dizer que essas exigências são para funções de onda normalizadas no sentido mais tradicional. É possível a existência de sistemas quânticos cuja função de onda não seja normalizável no sentido que estamos discutindo aqui, e que podem ser tratados por outras abordagens. Contudo, essas são exceções e exigem um conhecimento mais avançado. Neste momento, devemos nos ater aos casos que obedecem a essas condições, por serem tanto mais comuns quanto mais simples de trabalhar.

Mas uma vez que o termo entre parênteses vai então mais rápido a $0$ quando $x\rightarrow\infty$, o primeiro termo é $0$ e ficamos com:

\begin{equation}
\frac{d\left\langle x\right\rangle }{dt}=-\frac{i\hbar}{2m}\int\left(\Psi^{*}\frac{\partial\Psi}{\partial x}-\frac{\partial\Psi^{*}}{\partial x}\Psi\right)dx
\end{equation}

Ainda vale comentar que Griffiths define que os estados fisicamente possíveis correspondem a soluções quadrado-integráveis da equação de Schrödinger \footnote{$\int^{+\infty}_{-\infty}\left|f\left(x\right)\right|^{2}dx<\infty$(\href{https://en.wikipedia.org/wiki/Square-integrable_function}{definição})}.Estendendo essa discussão, considerando que o comportamento de $|f(x)|^{2}$ pode ser aproximado por $|x|^{-p}$ no infinito, temos então:

\begin{align}
\int^{+\infty}_{-\infty}\frac{dx}{\left|x\right|^{p}} & =\int^{+\infty}_{-\infty}\frac{dx}{\left|x\right|^{p}}+=\int^{0}_{-\infty}\frac{dx}{\left(-x\right)^{p}}+\int^{+\infty}_{0}\frac{dx}{x^{p}}\\
 & =-\int^{0}_{+\infty}\frac{du}{u^{p}}+\int^{+\infty}_{0}\frac{dx}{x^{p}}\\
 & =\int^{+\infty}_{0}\frac{du}{u^{p}}+\int^{+\infty}_{0}\frac{dx}{x^{p}}\\
 & =2\int^{+\infty}_{0}\frac{dx}{x^{p}}\\
 & =-\int^{+\infty}_{0}x^{-p}dx\\
 & =\left.\frac{x^{-p+1}}{p-1}\right|^{+\infty}_{0}\\
 & =\left(p-1\right)\lim_{x\rightarrow\infty}\frac{1}{x^{p-1}}-0\\
 & =\left(p-1\right)\lim_{x\rightarrow\infty}\frac{1}{x^{p-1}}
\end{align}

Para que a integral não diverja, precisamos que$p-1\geq0$, ou seja, $p\geq1,$logo, o menor valor que $p$ pode assumir para a integral convergir é $p=1$. Como consideramos que o comportamento assintótico é $|f(x)|^{2}\sim\frac{1}{|x|^{p}}$, temos $|f(x)|\sim\frac{1}{\sqrt{\left|x\right|^{p}}}$ como Griffiths menciona\footnote{Mais especificamente, se $p$ for real, temos $\left|x^{p}\right|$=$\left|x\right|^{p}$, logo $|f(x)|\sim\frac{1}{|\sqrt{\left|x^{p}\right|}}$.}. Pois:

\begin{equation}
\left|f\left(x\right)\right|^{2}=\left|\frac{1}{\sqrt{x^{p}}}\right|^{2}=\frac{1}{\left|x^{p}\right|}
\end{equation}

Como $p$ deve ser maior que 1 para a normalização da função de onda, o valor mínimo que $f(x)$ pode assumir para garantir isso é $f(x)\sim\frac{1}{\sqrt{|x|}}.$ Além disso, podemos notar que, conforme aumentamos $p$, a função $f(x)$ deve decair cada vez mais rápido para zero quando $x\to\infty$. É nesse sentido que se diz que, evidentemente, $\Psi(x,t)$ deve ir a zero mais rápido do que $\frac{1}{\sqrt{|x|}}$ quando $|x|\to\infty$.

Mas este caso, apesar de ser uma aproximação útil e cobrir a maioria dos casos fisicamente válidos, como foi discutido, não cobre necessariamente todos. É possível que uma função de onda fisicamente válida possua outra aproximação assintótica ou mesmo não seja quadrado-integrável. Neste último caso, é necessário usar alguma generalização da condição de normalização. Para a função de onda ser válida fisicamente, é necessário principalmente que ela seja uma solução da equação de Schrödinger e que possua uma interpretação probabilística (que pode ser feita através de uma generalização da normalização), além, evidentemente, de obedecer às condições impostas pelo problema físico.

Por enquanto, trabalharemos com os casos em que essas condições a respeito da normalização da função de onda são válidas, isto é, que a função seja quadrado-integrável e, consequentemente, apresente uma normalização do tipo$\int^{+\infty}_{-\infty}|\Psi(x,t)|^{2}\,dx=1,$ além de possuir um comportamento assintótico no infinito similar a$\frac{1}{|x|^{p}},$ com $p\geq\frac{1}{2}$\footnote{Vale destacar que, neste último cálculo do limite, precisamos apenas que o produto entre a função de onda e sua derivada decaia mais rapidamente que $\frac{1}{x}$ quando $x\to\infty$. Com a imposição adotada de que a função de onda decaia no infinito, no mínimo, como $\frac{1}{\sqrt{x}}$, garantimos essa condição. Por outro lado, se a função de onda decaísse, por exemplo, como $e^{-|x|}$, isso significaria que ela decai ainda mais rapidamente, e também teríamos o mesmo resultado no limite. No entanto, precisaríamos calcular para verificar se, com esse comportamento assintótico, a função de onda é quadrado-integrável.}.

Faremos isso da mesma forma que é assumido por Griffiths, pois assim podemos trabalhar com os casos mais gerais e deixar as exceções para um estudo mais especializado no futuro. Retomando, então, a discussão do limite fazendo uso dessas condições, ou seja, considerando que o comportamento assintótico da função de onda no infinito é similar a $\frac{1}{\sqrt{x}}$, temos: 
\begin{equation}
\Psi^{*}\frac{\partial\Psi}{\partial x}\sim\frac{1}{\sqrt{x}}\left(-\frac{1}{2x^{\frac{3}{2}}}\right)=-\frac{1}{2x^{2}}
\end{equation}

E consequentemente

\begin{equation}
\lim_{x\rightarrow\infty}x\left(\Psi^{*}\frac{\partial\Psi}{\partial x}-\frac{\partial\Psi^{*}}{\partial x}\Psi\right)\sim\lim_{x\rightarrow\infty}\left(-\frac{x}{2x^{2}}\right)=\lim_{x\rightarrow\infty}\left(-\frac{1}{2x}\right)=0
\end{equation}

Então ficamos com:

\begin{equation}
\frac{d\left\langle x\right\rangle }{dt}=-\frac{i\hbar}{2m}\int\left(\Psi^{*}\frac{\partial\Psi}{\partial x}-\frac{\partial\Psi^{*}}{\partial x}\Psi\right)dx
\end{equation}

Usando novamente a integral por substituição, para a segunda integral fazendo $f=\Psi$ e $g=\Psi^{*}$

\begin{align}
\int^{b}_{a}f\frac{dg}{dx}dx & =\left[fg\right]^{b}_{a}-\int^{b}_{a}\frac{df}{dx}gdx\\
\int^{\infty}_{-\infty}\frac{\partial\Psi^{*}}{\partial x}\Psi dx & =\left[\Psi\Psi^{*}\right]^{+\infty}_{-\infty}-\int^{\infty}_{-\infty}\Psi^{*}\frac{\partial\Psi}{\partial x}\\
 & =-\int^{\infty}_{-\infty}\Psi^{*}\frac{\partial\Psi}{\partial x}
\end{align}

Onde usamos o fato de que a função de onda vai a 0 no infinito . Temos então:

\begin{equation}
\frac{d\left\langle x\right\rangle }{dt}=-\frac{i\hbar}{2m}\int\left(\Psi^{*}\frac{\partial\Psi}{\partial x}+\Psi^{*}\frac{\partial\Psi}{\partial x}\right)dx=-\frac{i\hbar}{m}\int\Psi^{*}\frac{\partial\Psi}{\partial x}dx
\end{equation}

Vamos postular\footnote{Postulado é uma afirmação fundamental aceita como verdadeira sem necessidade de demonstração, que serve de base para o desenvolvimento de um sistema lógico ou teórico.}: O valor esperado da velocidade é igual a derivada no tempo do valor esperado da posição:

\begin{equation}
\left\langle v\right\rangle =\frac{d\left\langle x\right\rangle }{dt}
\end{equation}

Temos então o momento:

\begin{equation}
\left\langle p\right\rangle =m\frac{d\left\langle x\right\rangle }{dt}=-i\hbar\int\Psi^{*}\frac{\partial\Psi}{\partial x}dx
\end{equation}

Podemos reescrever as expressões como:

\begin{align}
\left\langle x\right\rangle  & =\int\Psi^{*}\left[x\right]\Psi dx\\
\left\langle p\right\rangle  & =\int\Psi^{*}\left[-i\hbar\left(\frac{\partial}{\partial x}\right)\right]\Psi dx
\end{align}

Dizemos que o operador \footnote{Um operador é uma entidade matemática que atua sobre o espaço de estados físicos de um sistema, isto é, o conjunto de todos os estados possíveis que o sistema pode assumir. Quando aplicado a um estado físico, o operador gera outro estado físico ou uma informação associada a esse estado, dependendo do tipo de operador.} $x$ “representa” a posição e o operador $-i\hbar\frac{\partial}{\partial x}$ “representa” o momento. Para calcular os valores esperados, fazemos um “sanduíche”, colocando o operador entre $\Psi^{*}$ e $\Psi$. Todas as variáveis dinâmicas clássicas podem ser expressas em termos de posição e momento. Por exemplo, a energia cinética é dada por:

\begin{equation}
T=\frac{1}{2}mv^{2}=\frac{p^{2}}{m}
\end{equation}

O momento angular é dado por:

\begin{equation}
\boldsymbol{L}=\boldsymbol{r}\times m\boldsymbol{v}=\boldsymbol{r}\times\boldsymbol{p}
\end{equation}

Então, para calcular o valor esperado de qualquer quantidade $Q(x,p)$, substituímos $p$ por $-i\hbar\frac{\partial}{\partial x}$ e inserimos o operador resultante entre o conjugado$\Psi^{*}$ e a própria função de onda $\Psi$. De forma geral:

\begin{equation}
\left\langle Q\left(x,p\right)\right\rangle =\int\Psi^{*}\left[Q\left(x,-i\hbar\frac{\partial}{\partial x}\right)\right]\Psi dx\label{eq:24}
\end{equation}

Griffiths promete fundamentar essa equação em uma base teórica sólida futuramente, mas por enquanto podemos considerá-la um axioma\footnote{Basicamente, é o mesmo que um postulado. A diferença entre axioma e postulado é mais uma questão de contexto: usa-se mais “axioma” em matemática e lógica, enquanto “postulado” é mais comum em outras ciências. Além disso, axioma costuma ser visto como algo mais universal, enquanto postulado está relacionado a uma ciência ou teoria específica.} . Por exemplo, o valor esperado para a energia cinética pode ser escrito como:

\begin{equation}
\left\langle T\right\rangle =-\frac{\hbar^{2}}{m}\int\Psi^{*}\frac{\partial^{2}\Psi}{\partial x^{2}}dx
\end{equation}

\textbf{Exemplo: Calcular $d\left\langle p\right\rangle /dt$}

Queremos obter:

\begin{equation}
\frac{d\left\langle p\right\rangle }{dt}=\left\langle -\frac{\partial V}{\partial x}\right\rangle 
\end{equation}

Vamos retomar a definição de $\left\langle p\right\rangle $

\begin{align}
\frac{d\left\langle p\right\rangle }{dt} & =\frac{d}{dt}\left(-i\hbar\int\Psi^{*}\frac{\partial\Psi}{\partial x}dx\right)\\
 & =-i\hbar\frac{d}{dt}\left(\int\Psi^{*}\frac{\partial\Psi}{\partial x}dx\right)\\
 & =-\left[\int\left(i\hbar\frac{\partial\Psi^{*}}{\partial t}\frac{\partial\Psi}{\partial x}+i\hbar\Psi^{*}\frac{\partial}{\partial t}\left(\frac{\partial\Psi}{\partial x}\right)\right)dx\right]\\
 & =\left[\int\left(\left(-i\hbar\frac{\partial\Psi^{*}}{\partial t}\right)\frac{\partial\Psi}{\partial x}-\Psi^{*}\frac{\partial}{\partial x}\left(i\hbar\frac{\partial\Psi}{\partial t}\right)\right)dx\right]
\end{align}

Onde usamos a simetria das segundas derivadas\footnote{$\frac{\partial}{\partial x}\left(\frac{\partial f}{\partial y}\right)=\frac{\partial}{\partial y}\left(\frac{\partial f}{\partial x}\right)$(\href{https://en.wikipedia.org/wiki/Symmetry_of_second_derivatives}{prova}).} e usando agora a equação de Schrödinger (e o complexo conjugado dela)

\begin{equation}
i\hbar\frac{\partial\Psi}{\partial t}=-\frac{\hbar^{2}}{2m}\frac{\partial^{2}\Psi}{\partial x^{2}}+V\Psi
\end{equation}

Temos:

\begin{align}
\frac{d\left\langle p\right\rangle }{dt} & =-\left[\int\left(i\hbar\frac{\partial\Psi^{*}}{\partial t}\frac{\partial\Psi}{\partial x}+i\hbar\Psi^{*}\frac{\partial}{\partial x}\left(\frac{\partial\Psi}{\partial t}\right)\right)dx\right]\\
 & =\int\left(\left(-\frac{\hbar^{2}}{2m}\frac{\partial^{2}\Psi^{*}}{\partial x^{2}}+V\Psi^{*}\right)\frac{\partial\Psi}{\partial x}-\Psi^{*}\frac{\partial}{\partial x}\left(-\frac{\hbar^{2}}{2m}\frac{\partial^{2}\Psi}{\partial x^{2}}+V\Psi\right)\right)dx\\
 & =\int\left(-\frac{\hbar^{2}}{2m}\frac{\partial^{2}\Psi^{*}}{\partial x^{2}}\frac{\partial\Psi}{\partial x}+V\Psi^{*}\frac{\partial\Psi}{\partial x}+\frac{\hbar^{2}}{2m}\Psi^{*}\frac{\partial^{3}\Psi}{\partial x^{3}}-\Psi^{*}\frac{\partial}{\partial x}\left(V\Psi\right)\right)dx\\
 & =\int\left[\frac{\hbar^{2}}{2m}\left(-\frac{\partial^{2}\Psi^{*}}{\partial x^{2}}\frac{\partial\Psi}{\partial x}+\Psi^{*}\frac{\partial^{3}\Psi}{\partial x^{3}}\right)+\left(V\Psi^{*}\frac{\partial\Psi}{\partial x}-\Psi^{*}\frac{\partial}{\partial x}\left(V\Psi\right)\right)\right]dx\\
 & =\int\left[\frac{\hbar^{2}}{2m}\left(-\frac{\partial^{2}\Psi^{*}}{\partial x^{2}}\frac{\partial\Psi}{\partial x}+\Psi^{*}\frac{\partial^{3}\Psi}{\partial x^{3}}\right)+\left(V\Psi^{*}\frac{\partial\Psi}{\partial x}-\Psi^{*}\frac{\partial V}{\partial x}\Psi-\Psi^{*}V\frac{\partial}{\partial x}\Psi\right)\right]dx\\
 & =\frac{\hbar^{2}}{2m}\int\left(-\frac{\partial^{2}\Psi^{*}}{\partial x^{2}}\frac{\partial\Psi}{\partial x}+\Psi^{*}\frac{\partial^{3}\Psi}{\partial x^{3}}\right)dx-\int\Psi^{*}\frac{\partial V}{\partial x}\Psi dx
\end{align}

Na primeira integral precisamos integrar por partes. Pegando a segunda integral fazendo $f=\Psi^{*}$ e $g=\partial^{2}\Psi/\partial x^{2}$ então:

\begin{align}
\int^{b}_{a}f\frac{dg}{dx}dx & =\left[fg\right]^{b}_{a}-\int^{b}_{a}\frac{df}{dx}gdx\\
\int\Psi^{*}\frac{\partial^{3}\Psi}{\partial x^{3}}dx & =\left[\Psi^{*}\frac{\partial^{2}\Psi}{\partial x^{2}}\right]^{+\infty}_{-\infty}-\int\frac{\partial\Psi^{*}}{\partial x}\frac{\partial^{2}\Psi}{\partial x^{2}}dx\\
 & =-\int\frac{\partial\Psi^{*}}{\partial x}\frac{\partial^{2}\Psi}{\partial x^{2}}dx
\end{align}

Então:

\begin{align}
\frac{d\left\langle p\right\rangle }{dt} & =\frac{\hbar^{2}}{2m}\int\left(-\frac{\partial^{2}\Psi^{*}}{\partial x^{2}}\frac{\partial\Psi}{\partial x}-\frac{\partial\Psi^{*}}{\partial x}\frac{\partial^{2}\Psi}{\partial x^{2}}\right)dx-\int\Psi^{*}\frac{\partial V}{\partial x}\Psi dx
\end{align}

Vamos aplicar a integral por partes novamente no segundo termo da primeira integral sendo agora $f=\partial\Psi^{*}/\partial x$ e $g=\partial\Psi/\partial x$

\begin{align}
\int^{b}_{a}f\frac{dg}{dx}dx & =\left[fg\right]^{b}_{a}-\int^{b}_{a}\frac{df}{dx}gdx\\
\int^{b}_{a}\frac{\partial\Psi^{*}}{\partial x}\frac{\partial\Psi^{2}}{\partial x^{2}}dx & =\left[\frac{\partial\Psi^{*}}{\partial x}\frac{\partial\Psi}{\partial x}\right]^{b}_{a}-\int^{b}_{a}\frac{\partial^{2}\Psi^{*}}{\partial x^{2}}\frac{\partial\Psi}{\partial x}dx\\
 & =-\int^{b}_{a}\frac{\partial^{2}\Psi^{*}}{\partial x^{2}}\frac{\partial\Psi}{\partial x}dx
\end{align}

Onde como no caso anterior, usamos a condição de que $\Psi$ e suas derivadas vão a zero no infinito. Temos então:

\begin{align}
\frac{d\left\langle p\right\rangle }{dt} & =\frac{\hbar^{2}}{2m}\int\left(-\frac{\partial^{2}\Psi^{*}}{\partial x^{2}}\frac{\partial\Psi}{\partial x}+\frac{\partial^{2}\Psi^{*}}{\partial x^{2}}\frac{\partial\Psi}{\partial x}dx\right)dx-\int\Psi^{*}\frac{\partial V}{\partial x}\Psi dx\\
 & =-\int\Psi^{*}\frac{\partial V}{\partial x}\Psi dxdx
\end{align}

E pela equação \ref{eq:24}, então:

\begin{align}
\left\langle Q\left(x,p\right)\right\rangle  & =\int\Psi^{*}\left[Q\left(x,-i\hbar\frac{\partial}{\partial x}\right)\right]\Psi dx\\
\left\langle -\frac{\partial V}{\partial x}\right\rangle  & =\int\Psi^{*}\left[-\frac{\partial V}{\partial x}\right]\Psi dx
\end{align}

Logo:

\begin{equation}
\frac{d\left\langle p\right\rangle }{dt}=\left\langle -\frac{\partial V}{\partial x}\right\rangle 
\end{equation}

Como queríamos demonstrar.

\subsection{A relação da incerteza}

Se pegarmos uma corda, temos duas situações extremas: 
\begin{itemize}
\item Se a oscilamos periodicamente, temos claramente um comprimento de onda, mas não podemos determinar onde ela está localizada. 
\item Se enviamos apenas um pulso, podemos localizar o pulso, mas não podemos falar em comprimento de onda. 
\end{itemize}
Casos intermediários apresentam uma situação de troca entre uma condição e outra. Isso se aplica a qualquer fenômeno ondulatório, inclusive à função de onda da mecânica quântica. Embora exista um teorema que pode ser demonstrado rigorosamente por meio da análise de Fourier, vamos nos concentrar aqui no argumento qualitativo. A fórmula de de Broglie\footnote{Muitos autores tomam a fórmula de de Broglie como um axioma, a partir do qual deduzem a associação do momento com o operador $-i\hbar\frac{\partial}{\partial x}$.} relaciona o momento da partícula à função de onda:

\begin{equation}
p=\frac{h}{\lambda}=\frac{2\pi\hbar}{\lambda}\label{eq:broglie}
\end{equation}

Temos, portanto, uma relação entre o momento e o comprimento de onda. Como essas grandezas são inversamente proporcionais, pode surgir a tendência de pensar que suas dispersões\footnote{Outro nome que a literatura usa para o desvio padrão $\sigma_{x}$é a dispersão, assim como se usa valor esperado para a média $\langle x\rangle$.} também sejam inversamente proporcionais --- mas isso é falso. A aproximação de primeira ordem para a propagação do erro de uma função $y=f(x)$, com uma única variável, é dada por(\href{https://web.chem.ox.ac.uk/teaching/Physics\%5C\%252520for\%5C\%252520CHemists/Errors/Propagation.html}{prova}):

\begin{equation}
\sigma_{y}\approx\left|\frac{df}{dx}\right|\sigma_{x}
\end{equation}

Ou seja:

\begin{equation}
\sigma_{p}\approx\left|-\frac{h}{\lambda^{2}}\right|\sigma_{\lambda}=\frac{h}{\lambda^{2}}\sigma_{\lambda}
\end{equation}

Ou ainda:

\begin{equation}
\frac{\sigma_{p}}{p}\approx\frac{\sigma_{\lambda}}{\lambda}
\end{equation}

Essa aproximação é suficiente para notarmos que chegamos a uma relação que expressa que a dispersão relativa de $p$é igual à dispersão relativa de $\lambda$. Assim, para um valor fixo de $\lambda$ (e, consequentemente, um valor fixo de \textbackslash ($p=h/\lambda$), um aumento na dispersão $\sigma_{\lambda}$ implica um aumento proporcional na dispersão $\sigma_{p}$, mesmo que os valores de $p$ e $\lambda$ sejam inversamente proporcionais. Quantitativamente o resultado que queremos é:

\begin{equation}
\sigma_{x}\sigma_{p}\geq\frac{\hbar}{2}
\end{equation}

Observação: A dispersão aqui se refere aos resultados obtidos em medições realizadas sobre sistemas identicamente preparados. Mesmo com a mesma preparação, os resultados não são necessariamente idênticos, mas se distribuem em torno de um valor médio.

\textbf{Exemplo: Uma partícula de massa $m$ com a seguinte função de onda}

\textbf{ 
\begin{equation}
\Psi\left(x,t\right)=A\exp\left[-a\left(\frac{mx^{2}}{\hbar}+it\right)\right]
\end{equation}
}

\textbf{Onde $A$ e $a$ são constantes reais positivas.}

\textbf{a) Queremos encontrar $A$}

\textbf{b) Encontrar qual função de energia potencial, $V\left(x\right)$ esta é uma solução para equação de Schrodinger}

\textbf{c) Calcular os valores esperados $x$, $x^{2}$, $p$ e $p^{2}$}

\textbf{d) Encontrar $\sigma_{x}$ e $\sigma_{p}$ e avaliar se o produto entre ambos é consistente com o princípio da incerteza}

Vamos começar calculando $A$, Primeiro:

\begin{align}
\left|\Psi\left(x,t\right)\right|^{2} & =\Psi^{*}\left(x,t\right)\Psi\left(x,t\right)\\
 & =\left(A\exp\left[-a\left(\frac{mx^{2}}{\hbar}+it\right)\right]\right)\left(A\exp\left[-a\left(\frac{mx^{2}}{\hbar}+it\right)\right]\right)\\
 & =A^{2}\exp\left[-\frac{2am}{\hbar}x^{2}\right]
\end{align}

Vamos substituir $u=\sqrt{\frac{2am}{\hbar}}x$ e $du=\sqrt{\frac{2am}{\hbar}}dx$ Integrando então:

\begin{align}
\int\left|\Psi\left(x,t\right)\right|^{2}dx & =\int A^{2}\exp\left[-\frac{2amx^{2}}{\hbar}\right]dx\\
1 & =A^{2}\sqrt{\frac{\hbar}{2am}}\int e^{-u^{2}}du
\end{align}

Essa é uma integral gaussiana de resultado bastante conhecido, mas vamos resolver uma vez (\href{https://github.com/jonasmaziero/mecanica_quantica/blob/main/notas_de_aula/02_Probabilidades.ipynb}{solução}) pelo menos. Definimos então uma integral:

\begin{equation}
G=\int^{+\infty}_{-\infty}e^{-ax^{2}}dx
\end{equation}

Então:

\begin{equation}
G^{2}=\left(\int^{+\infty}_{-\infty}e^{-ax^{2}}dx\right)^{2}=\left(\int^{+\infty}_{-\infty}e^{-ax^{2}}dx\right)\left(\int^{+\infty}_{-\infty}e^{-ax^{2}}dx\right)
\end{equation}

Ou como $x$ é a variável de integração em cada uma das integrais, ou seja, apenas um rótulo, podemos substituir $x\rightarrow y$ na segunda, e ficamos então com:

\begin{equation}
G^{2}=\left(\int^{+\infty}_{-\infty}e^{-ax^{2}}dx\right)\left(\int^{+\infty}_{-\infty}e^{-ay^{2}}dx\right)=\left(\int^{+\infty}_{-\infty}\int^{+\infty}_{-\infty}e^{-a\left(x^{2}+y^{2}\right)}dxdy\right)
\end{equation}

Agora podemos fazer uma mudança no sistema de coordenadas, obviamente, mudando também o elemento de área. Usando o sistema de coordenadas polares (\href{https://en.wikipedia.org/wiki/Polar_coordinate_system}{definição}):

\begin{align}
x & =r\cos\theta\\
y & =r\sin\theta\\
da & =rdrd\theta
\end{align}

Para cobrir todo o espaço, nossos limites são $0\leq r<\infty$ e $0\leq\theta<2\pi$.Substituindo:

\begin{align}
G^{2} & =\left(\int^{2\pi}_{0}\int^{\infty}_{0}e^{-ar^{2}}rdrd\theta\right)\\
 & =\left(\int^{2\pi}_{0}d\theta\right)\left(\int^{\infty}_{0}e^{-ar^{2}}rdr\right)\\
 & =2\pi\left(\int^{\infty}_{0}e^{-ar^{2}}rdr\right)
\end{align}

Substituindo $u=-ar^{2}$ então $\frac{du}{dr}=-a2r\rightarrow-\frac{du}{a}=2rdr$ e os limites de integração agora são $u=-a0^{2}=0$ e $u=-ar\infty^{2}=-\infty$

\begin{align}
G^{2} & =\pi\left(\int^{-\infty}_{0}e^{-ar^{2}}2rdr\right)\\
 & =-\frac{\pi}{a}\left(\int^{-\infty}_{0}e^{u}du\right)\\
 & =-\frac{\pi}{a}\left[e^{u}\right]^{-\infty}_{0}\\
 & =-\frac{\pi}{a}\left[\frac{1}{e^{u}}-e^{0}\right]^{-\infty}_{0}\\
 & =\frac{\pi}{a}
\end{align}

Logo se $G^{2}=\pi/a$ então $G=\sqrt{\pi/a}$ e:

\begin{equation}
\int^{+\infty}_{-\infty}e^{-ax^{2}}dx=G=\sqrt{\frac{\pi}{a}}
\end{equation}

Com isso podemos retomar e resolver nossa integral:

\begin{align}
1 & =A^{2}\sqrt{\frac{\hbar}{2am}}\int e^{-u^{2}}du\\
1 & =A^{2}\sqrt{\frac{\hbar}{2am}}\sqrt{\pi}\\
\sqrt{\frac{2am}{\hbar\pi}} & =A^{2}
\end{align}

Ou então:

\begin{equation}
A=\left(\frac{2am}{\hbar\pi}\right)^{\frac{1}{4}}
\end{equation}

Agora queremos saber para qual potencial de energia $V\left(x\right)$ essa é uma solução pra equação de Schrödinger. Partindo da equação de Schrödinger:

\begin{equation}
i\hbar\frac{\partial\Psi}{\partial t}=-\frac{\hbar^{2}}{2m}\frac{\partial^{2}\Psi}{\partial x^{2}}+V\Psi
\end{equation}

E derivando a função de onda:

\begin{align}
\Psi & =A\exp\left[-a\left(\frac{mx^{2}}{\hbar}+it\right)\right]\\
\frac{\partial\Psi}{\partial t} & =A\exp\left[-\frac{amx^{2}}{\hbar}\right]\left(-ai\right)\exp\left[-ait\right]=-ai\Psi\\
\frac{\partial\Psi}{\partial x} & =A\left[\left(-\frac{2amx}{\hbar}\right)\right]\exp\left[-\frac{amx^{2}}{\hbar}\right]\exp\left[-ait\right]=-\frac{2amx}{\hbar}\Psi\\
\frac{\partial^{2}\Psi}{\partial x^{2}} & =\frac{\partial}{\partial x}\left(-\frac{2amx}{\hbar}\Psi\right)=-\frac{2am}{\hbar}\Psi+\left(\frac{2amx}{\hbar}\right)^{2}\Psi
\end{align}

Então:

\begin{align}
i\hbar\frac{\partial\Psi}{\partial t} & =-\frac{\hbar^{2}}{2m}\frac{\partial^{2}\Psi}{\partial x^{2}}+V\Psi\\
V\Psi & =\frac{\hbar^{2}}{2m}\frac{\partial^{2}\Psi}{\partial x^{2}}+i\hbar\frac{\partial\Psi}{\partial t}\\
 & =\frac{\hbar^{2}}{2m}\left[-\frac{2am}{\hbar}\Psi+\left(\frac{2amx}{\hbar}\right)^{2}\Psi\right]+i\hbar\left[-ai\Psi\right]\\
 & =-a\hbar\Psi+2m\left(ax\right)^{2}\Psi+\hbar a\Psi\\
 & =\left[-a\hbar+2m\left(ax\right)^{2}+\hbar a\right]\Psi\\
\left[V\right]\Psi & =\left[2m\left(ax\right)^{2}\right]\Psi
\end{align}

Então:

\begin{equation}
V=2ma^{2}x^{2}
\end{equation}

Agora queremos calcular diversos valores esperados, vamos relembrar que já calculamos $\left|\Psi\left(x,t\right)\right|^{2}=A^{2}\exp\left[-\frac{2am}{\hbar}x^{2}\right]$. Começando com $\left\langle x\right\rangle $:

\begin{equation}
\left\langle x\right\rangle =\int xA^{2}\exp\left[-\frac{2am}{\hbar}x^{2}\right]dx
\end{equation}

Nesse caso, como o integrando é impar, isso é, se $G\left(x\right)=xA^{2}\exp\left[-\frac{2am}{\hbar}x^{2}\right]$, então 
\begin{align}
G\left(-x\right) & =-G\left(x\right)\\
\left(-x\right)A^{2}\exp\left[-\frac{2am}{\hbar}\left(-x\right)^{2}\right] & =-\left(xA^{2}\exp\left[-\frac{2am}{\hbar}x^{2}\right]\right)\\
-xA^{2}\exp\left[-\frac{2am}{\hbar}x^{2}\right] & =-xA^{2}\exp\left[-\frac{2am}{\hbar}x^{2}\right]
\end{align}

Então simplesmente a integral é zero e temos $\left\langle x\right\rangle =0$. Consequentemente se lembrarmos que $\left\langle p\right\rangle =m\frac{d\left\langle x\right\rangle }{dt}$, então $\left\langle p\right\rangle =0$. Agora vamos calcular $\left\langle x^{2}\right\rangle $:

\begin{equation}
\left\langle x^{2}\right\rangle =A^{2}\int x^{2}\exp\left[-\frac{2am}{\hbar}x^{2}\right]dx
\end{equation}

Escrevendo $\alpha=\frac{2am}{\hbar}$ então queremos resolver a integral:

\begin{equation}
I=\int x^{2}e^{-\alpha x^{2}}dx
\end{equation}

Fazendo integral por partes sendo $f=x$ e $g'=xe^{-\alpha x^{2}}$ precisamos resolver essa

\begin{equation}
g=\int xe^{-\alpha x^{2}}dx
\end{equation}

Substituindo $u=-\alpha x^{2}$ então $-\frac{du}{2\alpha}=-xdx$

\begin{equation}
g=-\frac{1}{2\alpha}\int e^{u}du=-\frac{e^{u}}{2\alpha}=-\frac{e^{-\alpha x^{2}}}{2\alpha}
\end{equation}

Evidentemente temos $f'=1$ Então:

\begin{align}
\int^{b}_{a}f\frac{dg}{dx}dx & =\left[fg\right]^{b}_{a}-\int^{b}_{a}\frac{df}{dx}gdx\\
\int x\left(xe^{-\alpha x^{2}}\right)dx & =\left[-\frac{xe^{-\alpha x^{2}}}{2\alpha}\right]^{+\infty}_{-\infty}+\frac{1}{2\alpha}\int e^{-\alpha x^{2}}dx
\end{align}

O primeiro termo vai a zero pois $e^{-x^{2}}$vai a zero ao infinito mais rapidamente que $\frac{1}{\left|x\right|}$ e a segunda integral é a gaussiana que a pouco vimos o resultado, então, lembrando que $\alpha=\frac{2am}{\hbar}$:

\begin{align}
\int x\left(xe^{-\alpha x^{2}}\right)dx & =\frac{1}{2\alpha}\sqrt{\frac{\pi}{\alpha}}=\frac{\sqrt{\pi}}{2}\left(\frac{\hbar}{2am}\right)^{\frac{3}{2}}
\end{align}

E sendo $A=\left(\frac{2am}{\hbar\pi}\right)^{\frac{1}{4}}$ então:

\begin{equation}
\left\langle x^{2}\right\rangle =A^{2}\int x^{2}\exp\left[-\frac{2am}{\hbar}x^{2}\right]dx=\left(\frac{2am}{\hbar\pi}\right)^{\frac{2}{4}}\frac{\sqrt{\pi}}{2}\left(\frac{\hbar}{2am}\right)^{\frac{3}{2}}
\end{equation}

Ou apenas:

\begin{equation}
\left\langle x^{2}\right\rangle =\frac{\hbar}{4am}
\end{equation}

E por fim para calcular $\left\langle p^{2}\right\rangle $ precisamos do operador $p$ dado por $-i\hbar\frac{\partial}{\partial x}$:

\begin{align}
\left\langle p^{2}\right\rangle  & =\int\Psi^{*}\left(-i\hbar\frac{\partial}{\partial x}\right)^{2}\Psi dx
\end{align}

Já calculamos anteriormente:

\begin{equation}
\frac{\partial^{2}\Psi}{\partial x^{2}}=-\frac{2am}{\hbar}\Psi+\left(\frac{2amx}{\hbar}\right)^{2}\Psi
\end{equation}

Então:

\begin{align}
\left\langle p^{2}\right\rangle  & =\int\Psi^{*}\left(-i\hbar\right)^{2}\left[-\frac{2am}{\hbar}\Psi+\left(\frac{2amx}{\hbar}\right)^{2}\Psi\right]\Psi dx\\
 & =\int\Psi^{*}\left(-\hbar^{2}\right)\left[-\frac{2am}{\hbar}\Psi+\left(\frac{2amx}{\hbar}\right)^{2}\Psi\right]\Psi dx\\
 & =\left(-\hbar^{2}\right)\left(-\frac{2am}{\hbar}\int\Psi^{*}\Psi dx+\left(\frac{2am}{\hbar}\right)^{2}\int x^{2}\Psi^{2}\Psi dx\right)\\
 & =\hbar2am\int\Psi^{*}\Psi dx-\left(2am\right)^{2}\int x^{2}\Psi^{2}\Psi dx
\end{align}

Pela normalização sabemos que a primeira integral vale $1$ e a segunda $\left\langle x^{2}\right\rangle =\frac{\hbar}{4am}$. Então:

\begin{align}
\left\langle p^{2}\right\rangle  & =2\hbar am-\left(2am\right)^{2}\frac{\hbar}{4am}\\
 & =2\hbar am-\hbar am
\end{align}

Então $\left\langle p^{2}\right\rangle =\hbar am$. Agora queremos a dispersão da posição e do momento, usando os valores médios calculados anteriormente:

\begin{align}
\sigma_{x} & =\sqrt{\left\langle x^{2}\right\rangle -\left\langle x\right\rangle ^{2}} & \sigma_{p} & =\sqrt{\left\langle p^{2}\right\rangle -\left\langle p\right\rangle ^{2}}\\
 & =\sqrt{\frac{\hbar}{4am}-0} &  & =\sqrt{\hbar am-0}\\
 & =\sqrt{\frac{\hbar}{4am}} &  & =\sqrt{\hbar am}
\end{align}

Então:

\begin{equation}
\sigma_{x}\sigma_{p}=\sqrt{\frac{\hbar}{4am}}\sqrt{\hbar am}=\frac{\hbar}{2}
\end{equation}

Este resultado é consistente com a relação de incerteza:

\begin{equation}
\sigma_{x}\sigma_{p}\geq\frac{\hbar}{2}
\end{equation}

\end{document}
